\documentclass[12pt]{article}
\usepackage[utf8]{inputenc}
\usepackage[T1]{fontenc}
\usepackage{hyperref}
\title{Krikati-Timbira Language Resources for K-12 Curriculum Development}
\author{Compiled by Claude Sonnet 5.5 (claude-sonnet-5-5) for Andr\'as Kornai}
\date{}
\begin{document}
\maketitle

\section{Introduction}
Krikati-Timbira (ISO 639-3: xri) is a Timbira language of the Northern Je (Jean) family, spoken in the Terra Ind\'igena Krikati in Maranh\~ao, Brazil. It is very closely related to Pykobj\^e; the main difference reported is phonological. This report is based on a small number of web searches. Little in the way of native-language classroom material turned up, and nothing below should be read as a complete survey.

\section{Primary Educational Resources}
No primer or textbook in Krikati was found in the search results. Research on Krikati bilingual schooling and teacher training exists (Portuguese-language academic work):
\begin{itemize}
\item Formation of Krikati indigenous teachers and bilingual pedagogical practice (university repository record): \url{https://univates.br/bdu/handle/10737/1053}
\item Study of the Timbira ethno-educational territory and Krikati protagonism in indigenous school education: \url{https://www.univates.br/bdu/handle/10737/2620}
\item Article in \emph{Revista Cocar} on the Krikati perspective on the Timbira ethno-educational territory: \url{https://periodicos.uepa.br/index.php/cocar/article/download/4343/1947/14489}
\end{itemize}
The results indicate that Krikati teachers have sought training in writing their mother tongue and in using it for literacy teaching.

\section{Linguistic and Reference Materials}
\begin{itemize}
\item Wikipedia entry on the Kr\~ikat\'i dialect: \url{https://en.wikipedia.org/wiki/Kr%C4%A9kat%C3%AD_dialect}
\item Wikipedia entry on the closely related Pykobj\^e dialect: \url{https://en.wikipedia.org/wiki/Pykobj%C3%AA_dialect}
\item Endangered Languages Project bibliography for the language: \url{https://endangeredlanguages.com/lang/2671/bibliography?hl=en}
\item Krikati number-word page (Max Planck EVA numerals database): \url{https://lingweb.eva.mpg.de/channumerals/Krikati-Timbira.htm}
\end{itemize}
A Krikati-Portuguese dictionary by a missionary author is cited in secondary sources (given as Pries 2008, \emph{Dicion\'ario Gavi\~ao-Krikati}); no URL for it was seen, and no full grammar was found.

\section{Religious and Community Resources}
No scripture translation or community organisation page was found in the search results.

\section{Digital and Audio Resources}
The Glosbe pages for the language code are crowdsourced and of uncertain quality: \url{https://glosbe.com/xri/lib}. No audio or video resources were found.

\section{Curriculum Development Context}
Documentation is thin for curriculum purposes. A curriculum would have to rely on work with Krikati teachers and community members, and on the Portuguese-language studies of Krikati bilingual education cited above. Further searching in Portuguese and consultation with the community are needed.

\end{document}